% Generated from the matching Typst scaffold by
% scripts/migrate_typst_report_to_latex.py.

% ── 1.2  Problem Statement ──────────────────────────────────────────────────
% PURPOSE: Name the pain, its consequence, and the unmet research gap.
% MOVE: Pain → Consequence → Gap ("therefore there exists a need …").
% BEATS: pain (bloat vs lossy pruning) · consequence (cost + entanglement) · gap.
% FIGURES/TABLES: none.
% CITATIONS: xiong2025 (more memory can hurt — experience-following),
%            alqithami2025 (closest six-policy prior — NOVELTY-THREAT).
% ─────────────────────────────────────────────────────────────────────────

\section{Problem Statement}

\begin{todobox}
Expand the beats below using Pain \textrightarrow{} Consequence \textrightarrow{} Gap, ending on an explicit "therefore there exists a need …". Pre-empt Alqithami (alqithami2025) as a novelty threat: conversational domain, retrieval not isolated.
\end{todobox}

% - PAIN: full memory bloats the context window; but naive pruning may drop
%   task-critical knowledge — both failure modes are real.
% - more memory can actively hurt via experience-following / interference,
%   not merely cost @xiong2025.
% - CONSEQUENCE: cost and latency rise with accumulation; worse, prior work
%   *entangles* the storage (forgetting) policy with the retrieval policy, so
%   no benefit is cleanly attributable — and none is demonstrated on coding.
% - NOVELTY-THREAT: @alqithami2025 is the closest six-policy comparison, but it
%   is conversational and does *not* hold retrieval constant — differentiate now.
% - GAP (therefore): there exists no controlled, *retrieval-fixed* comparison of
%   forgetting policies for coding agents — the central control this thesis adds.
